% Draft subsection for Section V (dataset audit). Numbers marked [RERUN] are
% pending the CSE-CIC-IDS2018 / Edge-IIoTset re-runs.
% Sources: results/diagnostics/leakage_audit/{README.md,schema.csv,features.csv,
% edgeiiot_layout.json,edgeiiot_repair_verification.{json,csv},layout_detector_control.json}
% Scripts: scripts/leakage_feature_audit.py, scripts/feature_names.py,
% scripts/verify_edgeiiot_repair.py, src/data/repair.py, tests/test_repair.py
% Generated 2026-09-17.
% LABELS TO VERIFY against the manuscript: tab:datasets (Table II; may be
% tab:datasets_overview), subsec:preprocessing (Sec. V-B has no visible label in
% the copy available here), and sec:results (the Section VI quantization
% analysis). tab:architecture_iot is confirmed correct.
% The ISCX-IDS2012 paragraph is repeated as a threat to validity at the end of
% this file, for Section VII.

\subsection{Feature-Level Leakage Audit}
\label{subsec:audit}

Published intrusion-detection corpora are distributed as derived feature tables,
and a defect in that derivation can encode the label in an input. We therefore
audit all eleven benchmarks before reporting any result, in two layers, and
release the audit scripts with the artifact.

\textbf{Layer~A (schema).} For every loader that concatenates several source
files we compare the column sets of those files. A column absent from some files
is padded on concatenation, so its ``is padded'' pattern is a source-file
indicator. One dataset fails: \mbox{CSE-CIC-IDS2018}, where
\texttt{02-20-2018.csv} alone carries \texttt{Flow~ID}, \texttt{Src~IP},
\texttt{Src~Port} and \texttt{Dst~IP}. The three string columns become
all-zero; \texttt{Src~Port} is non-zero on exactly the rows of that capture day,
48.1\% of the corpus, and two of the five classes (\emph{r2l}, \emph{u2r}) never
occur on that day. A single feature therefore rules both of them out on every row
where it fires. The GBDT teacher ranks it the most important
feature of the dataset, and it is the same feature that carries the first-layer
weight outlier analysed in Section~\ref{sec:results}: under per-tensor
quantization that outlier sets the scale for its whole layer, which is what made
the dense twin unstable on this corpus. The leak is therefore not merely
theoretical -- it is visible in both the teacher's attributions and the
student's weights. All four
columns are dropped, which also makes the pipeline consistent with the
identifier removal stated in Section~\ref{subsec:preprocessing}.

\textbf{Layer~B (preprocessed arrays).} On each training split we test, for every
feature, the indicator ``value above its floor'' against the label: its
normalised mutual information, and whether -- while the indicator is active on
between 5\% and 95\% of rows -- it is never (or always) active for some class,
i.e.\ whether observing it rules that class in or out on the rows where it fires. Flags are reviewed by name, since legitimate protocol features
also switch off for whole classes (r2l over TCP, ICMP floods without TCP flags).
Two datasets need action.

\emph{Edge-IIoTset} is derived from \texttt{tshark} output split on every comma
without quoting. Two consequences follow. First, the zero placeholder is written
\texttt{0} or \texttt{0.0} depending on the source capture: in
\texttt{mqtt.topic}, \texttt{mqtt.protoname}, \texttt{mqtt.conack.flags} and
\texttt{dns.qry.name.len} the spelling alone separates \emph{Normal} from every
attack class, and \texttt{http.request.version} / \texttt{http.request.method}
split the attacks into two fixed groups. Because the loader one-hot encodes these
columns as text, the spelling becomes a feature. Its normalised mutual
information with the label reaches 0.48, against 0.25 for the most informative
feature of the same corpus that the defect does not touch (\texttt{tcp.flags});
the decisive property, however, is not the magnitude but that the mapping is
deterministic.
Second, fields that \texttt{tshark} prints with an embedded comma shift the rest
of the row: \texttt{frame.time} (``Nov~24, 2021 \ldots'') and multi-valued fields
(\texttt{udp.port} is always \texttt{src,dst}; ICMP errors repeat
\texttt{ip.src\_host} and \texttt{ip.dst\_host} for the inner header). Every
\emph{DDoS\_UDP} row (121\,568) and every \emph{MITM} row (1\,214) is misaligned,
as is every benign UDP row -- 2\,227 rows, 0.14\% of the \emph{Normal} class, of
which the 1\,845 DNS rows are those the class-exclusion test above detects, the
remaining 382 carrying NTP, NetBIOS, mDNS or DHCP. The other twelve classes are
clean. The
same misalignment is present in the per-attack CSVs the dataset ships, so it
originates in the published extraction rather than in our merge.

\emph{ISCX-IDS2012} retains \texttt{direction}, a locality attribute derived from
the addresses: local-to-remote flows are 99.98\% benign and roughly 99\% of
attacks are local-to-local or remote-to-local by testbed construction. This is a
documented field of the corpus, not a defect of the derivation, so we keep it
rather than editing the benchmark, and report it here because it makes
ISCX-IDS2012 optimistic for every model alike.

\subsection{Repair and Verification}
\label{subsec:repair}

The \mbox{CSE-CIC-IDS2018} repair is the column drop above. For Edge-IIoTset we
canonicalise the zero spellings and hex renderings before encoding, drop
\texttt{http.referer} (payload text, including a Shellshock string on 185 rows)
and \texttt{dns.qry.name} (a hostname once the rows are aligned), and decode, by
counting the commas each field contributes, the three misaligned layouts we can
repair (\emph{MITM} being removed rather than decoded): attacker flood packets,
the victim's ICMP port-unreachable replies, and benign UDP. \emph{MITM} is garbled beyond a verifiable decode and is removed, leaving
fourteen classes. Benign UDP rows are decoded rather than dropped: removing them
would leave the repaired flood as the only UDP traffic in the corpus and create a
fresh shortcut. Each row carries a provenance flag (\emph{original} or which
decode was applied), released with the processed data so any result can be
recomputed with repaired rows excluded.

Every decoded row is verified against the packet captures the dataset ships: we
re-extract the same 61 fields with \texttt{tshark}~4.2.2, match each row to its
packet by (\texttt{ip.src\_host}, \texttt{ip.dst\_host}, \texttt{udp.port},
\texttt{udp.stream}) and compare field by field. All 123\,795 decoded rows match,
and all 60 comparable fields -- every field except the unrecoverable
\texttt{frame.time} -- agree with the capture on 100\% of rows in every layout.
The published values of the fields the decode moves agree on 0\% of rows: four
fields for the flood layout, seven for the ICMP-reply layout, and two
(\texttt{udp.stream}, \texttt{udp.time\_delta}) for benign UDP. Two details: benign UDP rows are
matched at flow level (a median of four packets share a flow key), so a row is
verified when some packet of its flow matches it on every field; and
\texttt{tshark}~4.x prints boolean fields as \texttt{True}/\texttt{False} where
the published extraction wrote \texttt{1}/\texttt{0}, which accounted for 59
apparent mismatches before the comparison was made version-agnostic. One field is
not repairable: \texttt{frame.time} is absent from the source CSV for every
\emph{DDoS\_UDP} row. It is not a model input, but the repair is not total and we
do not claim it is. Trailing Modbus fields truncated upstream are filled with
zero and agree with the captures on every row.

We also asked whether the misalignment is what makes \emph{DDoS\_UDP}
recognisable. A depth-four tree trained only on per-column type indicators
identifies the class with $F_1=1.000$ on the published data, but still reaches
0.999 after repair, while the same probe on a correctly extracted UDP flood
(CICIoT2023) reaches 0.701. The probe therefore does not isolate the artifact --
in Edge-IIoTset \emph{DDoS\_UDP} is the only attack class carrying UDP fields --
and we report the field-level verification above instead.

The Edge-IIoTset input width changes from 95 to 64 for three separable reasons:
collapsing the two zero spellings and parsing hex removes 18 one-hot columns, the
two dropped columns remove 12, and removing \emph{MITM} removes one level that
only that class exhibited; realignment is width-neutral, since the hostnames it
moves out of \texttt{dns.qry.name.len} leave with \texttt{dns.qry.name}.
Deduplicated row count changes from 1\,909\,671 to 1\,909\,313, and the class
count from 15 to 14 -- which also moves the macro-$F_1$ denominator, the output
dimension in Table~\ref{tab:architecture_iot} and the taxonomy count in
Table~\ref{tab:datasets}, independently of the leakage repair itself.

All affected results were regenerated from the repaired corpora. [RERUN: state
the direction and magnitude plainly, e.g. ``macro-$F_1$ falls by X.X points on
CSE-CIC-IDS2018 and Y.Y on Edge-IIoTset, for every model''; do not assert a
direction before the runs report it.] The repaired figures are the ones reported
throughout this paper; the pre-repair figures are recoverable from the released
provenance flags.

% ---------------------------------------------------------------------------
% For Section VII (threats to validity): the one audit finding we disclose
% rather than repair.
\paragraph{Residual dataset shortcut.}
ISCX-IDS2012 retains \texttt{direction}, an address-derived locality attribute:
local-to-remote flows are 99.98\% benign, and roughly 99\% of attacks are
local-to-local or remote-to-local because of how the testbed staged them. We keep
the field, since removing it would edit the benchmark and make our numbers
incomparable with published work, but results on this corpus are optimistic for
every model alike and the dataset contributes correspondingly little
discrimination between them.
